\documentclass[11pt,a4paper]{article}
\usepackage[T1]{fontenc}
\usepackage[margin=1in]{geometry}
\usepackage{graphicx}
\usepackage{amsmath}
\usepackage{amssymb}
\usepackage[authoryear,round]{natbib}
\usepackage{hyperref}
\usepackage{booktabs}
\usepackage{float}
\usepackage{xcolor}
\usepackage{textcomp}
\usepackage{lineno}
\usepackage{setspace}
\usepackage{enumitem}
\usepackage{tabularx}
\usepackage{multirow}
\usepackage{longtable}

\linenumbers
\onehalfspacing

% Title for npj Urban Sustainability (must be <=15 words, NO punctuation -- see reframe note; current title is a PLACEHOLDER pending reframe)
\title{\Large\textbf{School-age PM2.5 exposure and classroom protection across six Central Asian capitals}}

\author{Abduxoliq Ashuraliyev$^{1,*}$}

\date{}

\begin{document}

\maketitle

\vspace{-0.5cm}
\noindent\rule{\textwidth}{0.5pt}

\begin{center}
$^{1}$Independent Researcher, Open Data Initiative, Tashkent, Uzbekistan

$^{*}$Corresponding author: Abduxoliq Ashuraliyev (jack00040008@outlook.com)
\end{center}

\noindent\vspace{0.3cm}

% Abstract (npj Urban Sustainability: <=150 words, NO subheadings -- currently ~165 words, needs trimming in format pass)
\begin{abstract}
\noindent Central Asian cities are emerging air-pollution hotspots, yet children's exposure and protection remain unquantified. Using reference-grade monitors at six U.S. embassies, we characterize school-age PM2.5 exposure across the region's capitals (January 2022--June 2023) and demonstrate a low-resource framework transferable to data-scarce cities. Five of six capitals exceeded the WHO annual guideline 4.6--10.7-fold, with school-hours concentrations of 17.5--52.1~µg/m$^3$ and steep winter amplification from heating and inversions; Dushanbe was worst (annual 53.3~µg/m$^3$) and Tashkent mid-range (37.9). For Tashkent we map 605 schools: 15\% lie within 100~m of major roads, and modeled classroom concentrations span 18.9--33.9~µg/m$^3$---a within-city equity gradient by building age. Paediatric exposure--response functions attribute a large share of childhood asthma to this exposure. Classroom HEPA filtration---about 3\% of documented citywide health costs---offers near-term, equity-aware protection where ambient controls require decades, providing a template for data-scarce cities region-wide.
\end{abstract}

\noindent\textbf{Keywords:} urban air quality; school siting; PM2.5; built environment; traffic corridors; ventilation policy; Central Asia; sustainable cities

\vspace{0.5cm}
\noindent\rule{\textwidth}{0.5pt}

%----------------------------------------------------------------------
\section{Introduction}
%----------------------------------------------------------------------

Central Asian cities have emerged as global hotspots of air pollution, yet remain among the least-studied regions for air quality-health relationships~\cite{worldbank2024centralasia}. Recent assessments rank Central Asian capitals among the world's most polluted, with PM2.5 concentrations exceeding WHO guidelines by 4--13-fold despite relatively low population densities---indicating severe anthropogenic impact and limited technological penetration for emission control~\cite{tursumbayeva2023}. As documented by Kerimray et al. (2023), air quality research publications for Central Asian cities are "orders of magnitude lower" than for comparably polluted cities in other regions, creating a severe evidence-publication disparity~\cite{tursumbayeva2023}. This stark knowledge gap leaves policymakers in a region of 75 million people without the exposure assessments required for evidence-based urban planning and child health protection.

Independent assessments rank three Central Asian capitals---Dushanbe, Tashkent, and Bishkek---among the world's thirty most polluted, with documented air-pollution health costs ranging from 0.7\% of GDP in Uzbekistan to as much as 16.8\% of city-level gross regional product in Kazakhstan, depending on city and valuation method~\cite{worldbank2024centralasia,unece2024uzbekistan,tursumbayeva2023}. These assessments, however, characterize ambient, whole-population exposure; none resolves what children breathe during school hours. We therefore apply a single common framework---one reference-grade monitor per city combined with open spatial and health-burden data---across all six capitals, using Tashkent as the worked example for the spatial, infiltration, health, and intervention analyses that follow.

Children spend 30--40 hours weekly in schools during critical developmental years, making the school microenvironment a decisive determinant of cumulative exposure~\cite{gauderman2015association,gehring2013air}. Where schools sit near busy roads, near-road microenvironments add a traffic-related increment to the ambient burden~\cite{who2021guidelines,apte2017high}. Fine particulate matter (PM2.5)---the most health-relevant air pollutant---penetrates deep into the respiratory system, translocates to systemic circulation, and exhibits particularly severe effects on developing children due to higher ventilation rates per unit body weight and ongoing lung development~\cite{pope2004cardiovascular,brook2010particulate,schwartz2004air}.

Tashkent, Uzbekistan (population 2.6 million; 2.5 million metropolitan vehicles) exemplifies this urban sustainability challenge. World Bank analyses estimate that air pollution costs Uzbekistan 6.5\% of GDP annually, with Tashkent-specific health costs reaching \$488 million (0.7\% of national GDP)~\cite{worldbank2024uzbekistan,unece2024uzbekistan}. The city's enclosed basin geography, winter thermal inversions, and coal-dependent heating create seasonal pollution episodes with documented population-health consequences~\cite{ashuraliyev2025pm25tashkent}---yet Tashkent has no school ventilation standards, no systematic child-exposure assessment, and no policy framework for classroom intervention.

Critically, Tashkent's ambient PM2.5 now exceeds that of Beijing, which achieved substantial improvements through aggressive policy intervention (annual mean 29--33~µg/m$^3$ by 2022--2023 versus Tashkent's 37.9~µg/m$^3$)~\cite{crea2023china}. This comparison demonstrates both the severity of Central Asian exposure and the policy opportunity: school-based interventions proven effective in China and India offer transferable solutions before ambient pollution is addressed through long-term emission control~\cite{landrigan2017lancet,khreis2017exposure}.

Data-scarce LMICs rarely have dense monitoring networks for school-specific air quality assessments~\cite{cohen2017estimates}. This study develops a scalable framework requiring only: (1) a single reference-grade monitor (U.S.\ Embassy stations operate in most capitals), (2) published building infiltration literature, and (3) open-source health burden functions. We demonstrate that this minimal-data approach---replicable across LMIC cities facing similar constraints---can generate actionable policy evidence within available resource envelopes.

A 2024 World Bank assessment using chemical-transport modelling has since established the authoritative ambient and emissions baseline for Tashkent---a population-weighted annual mean of 38.8~µg/m$^3$, approximately 3,000 attributable premature deaths annually, and a source mix dominated by residential heating in winter (28\% of annual concentration, up to $\sim$45\% in peak months) and windblown dust (36\%) rather than transport (16\%)~\cite{worldbank2024tashkent}. That assessment treats the whole population and ambient-emissions policy; it does not resolve where children are exposed, how classroom environments modify that exposure, or which near-term protections are warranted. We address this gap. This study asks four questions relevant to urban sustainability policy: (i) what is the magnitude and temporal pattern of school-hour PM2.5 exposure relative to WHO guidelines; (ii) how are schools distributed relative to major-road microenvironments, and how does building-envelope quality create within-city exposure inequities; (iii) what share of childhood respiratory and asthma burden is attributable to this exposure; and (iv) how cost-effective is classroom filtration as a near-term, equity-aware protection. By building the analysis from a single open reference monitor combined with open spatial and health-burden data, we also demonstrate a minimal-resource template transferable to the many data-scarce cities that share Tashkent's constraints~\cite{worldbank2024dialogue}.

%----------------------------------------------------------------------
\section{Results}
%----------------------------------------------------------------------

\subsection{PM2.5 Concentrations Exceed Guidelines Year-Round}

Analysis of 8,301 hourly PM2.5 measurements from the U.S.\ Department of State's air quality monitoring station at the U.S.\ Embassy in Tashkent (January 2022--June 2023) revealed chronic guideline exceedances (Table~\ref{tab:descriptive}). Mean PM2.5 concentration was 37.9~µg/m$^3$ (annual average computed as the mean of 346 daily means; arithmetic mean of the 8,301 hourly values 37.8~µg/m$^3$, SD 38.4), with median 27.0~µg/m$^3$ and 75th percentile 42.0~µg/m$^3$. This mean represents 7.6-fold the WHO 2021 annual guideline (5~µg/m$^3$) and exceeds WHO Interim Target 2 (25~µg/m$^3$)~\cite{who2021guidelines}.

\begin{table}[H]
\centering
\caption{Descriptive statistics for PM2.5 concentrations (µg/m$^3$) in Tashkent, Uzbekistan (n=8,301 hourly measurements, January 2022--June 2023). Values describe the hourly distribution; the hourly arithmetic mean is 37.8~µg/m$^3$, while the WHO-comparison annual average (mean of 346 daily means) is 37.9~µg/m$^3$. Data from U.S.\ Embassy monitoring station via OpenAQ (Station ID 8881, StateAir program).}
\label{tab:descriptive}
\begin{tabular}{lcccccccc}
\toprule
& Mean & Median & SD & Min & Max & P25 & P75 & P95 \\
\midrule
PM2.5 & 37.8 & 27.0 & 38.4 & 2.0 & 857.0 & 19.0 & 42.0 & 108.0 \\
\bottomrule
\end{tabular}
\end{table}

The 24-hour WHO guideline (15~µg/m$^3$) was exceeded on 322 of 346 analyzed days (93.1\%). The maximum hourly value (857.0~µg/m$^3$) was observed during an extreme winter pollution episode; this value is within the BAM-1020 measurement range (0--1000~µg/m$^3$) and consistent with documented severe winter episodes in Central Asian cities during thermal inversions~\cite{worldbank2024centralasia,tursumbayeva2023}. Guideline exceedances were pervasive across all seasons (Supplementary Table S1).

\subsection{Seasonal Patterns Reflect Urban Heating and Atmospheric Dynamics}

PM2.5 concentrations demonstrated pronounced seasonal variation (Fig.~\ref{fig:seasonal}). Winter (December--February) mean concentrations reached 59.1~µg/m$^3$ (SD 53.2), with 98.2\% of winter days exceeding the 24-hour guideline. Summer (June--August) concentrations averaged 26.1~µg/m$^3$ (SD 13.2), representing a 126\% winter-to-summer increase (Kruskal-Wallis $H$=847.3, $p<$0.001). This pattern reflects thermal inversion trapping of heating emissions, reduced atmospheric ventilation, and cold-start vehicle emissions characteristic of continental climates~\cite{zhang2017drivers,li2017wintertime}. December showed the highest monthly mean (78.5~µg/m$^3$), coinciding with peak residential heating demand.

\begin{figure}[H]
\centering
\includegraphics[width=0.95\textwidth]{tashkent_daily_means.png}
\caption{Daily mean PM2.5 concentrations in Tashkent (January 2022--June 2023). Red dashed line indicates WHO 24-hour guideline (15~µg/m$^3$). Pronounced winter peaks coincide with heating season and atmospheric inversions. Of 346 analyzed days ($\geq$18h data completeness), 322 (93.1\%) exceeded the WHO 24-hour guideline, with winter months showing the most severe concentrations.}
\label{fig:seasonal}
\end{figure}

Spring (March--May) and autumn (September--November) showed intermediate concentrations (26.9 and 37.5~µg/m$^3$, respectively) with elevated variability reflecting transitional meteorology and episodic dust transport from surrounding deserts (Supplementary Fig.~S1).

\subsection{School-Hour and Commute Exposures Create Peak Child Burden}

Diurnal analysis revealed elevated PM2.5 during children's active hours (Fig.~\ref{fig:diurnal}). School hours (08:00--15:00 weekdays) averaged 34.1~µg/m$^3$, while morning commute periods (07:00--09:00) reached 39.8~µg/m$^3$---coinciding with peak traffic and atmospheric stability. Evening commute periods (16:00--18:00) showed secondary peaks of 31.8~µg/m$^3$. Nighttime minima (02:00--05:00) averaged 39.6~µg/m$^3$, elevated due to winter heating emissions, still exceeding WHO guidelines by nearly 8-fold.

\begin{figure}[H]
\centering
\includegraphics[width=0.95\textwidth]{tashkent_diurnal.png}
\caption{Diurnal PM2.5 patterns showing traffic-related morning (07:00--09:00) and evening (16:00--18:00) peaks. School hours (shaded) experience sustained elevated concentrations. Data represent mean hourly values across 18-month monitoring period (n=8,301 hourly measurements).}
\label{fig:diurnal}
\end{figure}

Weekday-weekend analysis revealed comparable school-hour concentrations, with weekday means (33.7~µg/m$^3$) slightly \textit{lower} than weekend means (35.1~µg/m$^3$). The near-absence of a weekday traffic signature---and overnight concentrations rivalling the commute peak---indicates that residential heating and regional dust, not local traffic, dominate the source mix, consistent with the World Bank's chemical-transport source apportionment~\cite{worldbank2024tashkent}. This source profile is consequential for policy: because the dominant drivers are diffuse and seasonal, the most tractable near-term lever for protecting children is the school microenvironment itself rather than localized traffic control~\cite{nieuwenhuijsen2019influence,morawska2019applications}.

\subsection{A Reference-Anchored School Exposure Surface}

To resolve \textit{where} children are exposed, we first established that the regional PM2.5 field is spatially near-uniform. Satellite-derived estimates (ACAG, 1~km) place the Tashkent annual mean within $\sim$10\% of our reference-monitor value, and the World Bank CAMx assessment reports inter-station correlation $r>0.85$ across the metropolitan area~\cite{worldbank2024tashkent,acag2022satpm25}. Coarse global reanalysis products do not resolve the city---a 25-point CAMS-global grid over the metropolitan area returns only two distinct values and under-reads the reference monitor by $\sim$50\%---underscoring that intra-urban differentials must be modelled from local predictors rather than read from global grids. We therefore treat the reference-monitor annual mean (37.9~µg/m$^3$) as a homogeneous background and model the within-city differential from two local predictors.

The first is near-road position. Geolocating 605 schools from OpenStreetMap and computing each school's distance to the nearest major road (motorway, trunk, or primary), we found the distribution strongly right-skewed: 93 schools (15.4\%) lie within 100~m of a major road and 153 (25.3\%) within 200~m. Applying a standard near-road decay (increment $= A\,e^{-d/L}$, with $A=5$~µg/m$^3$ at the roadside and $L=150$~m)~\cite{apte2017high,khreis2017exposure} yields outdoor exposures of 37.9--42.5~µg/m$^3$ across the network. The second is the building envelope; combining the outdoor surface with construction-era infiltration factors gives modelled indoor classroom concentrations of 19.0--33.9~µg/m$^3$ (Fig.~\ref{fig:exposuresurface}), a 1.7-fold gap between the most- and least-exposed deciles of schools. The building envelope dominates this differential while near-road position compounds it, together providing concrete criteria for school siting, renovation prioritisation, and dismissal-time traffic policy.

\begin{figure}[H]
\centering
\includegraphics[width=0.7\textwidth]{fig_school_exposure_surface.png}
\caption{Modelled indoor classroom PM2.5 across 605 Tashkent schools (OpenStreetMap), combining a satellite/CAMx-corroborated homogeneous regional background (37.9~µg/m$^3$), a near-road increment from road proximity (exponential decay, $A=5$~µg/m$^3$, $L=150$~m), and construction-era building-envelope infiltration. The reference monitor (Station 8881) is starred. Values are first-order proximity-model estimates pending in-classroom measurement.}
\label{fig:exposuresurface}
\end{figure}

\subsection{Indoor Infiltration Reveals Classroom Exposure Heterogeneity}

Building infiltration modeling using literature-derived indoor/outdoor ratios (0.50--0.80) appropriate for Central Asian building stock~\cite{abt2000relative,macintosh2009residential,chen2012review} estimated classroom PM2.5 concentrations spanning 18.9--30.3~µg/m$^3$ depending on building envelope quality (Table~\ref{tab:infiltration}).

\begin{table}[H]
\centering
\caption{Estimated school indoor PM2.5 concentrations by building sealing quality. Infiltration ratios based on published Central Asian building stock literature.}
\label{tab:infiltration}
\begin{tabular}{lccc}
\toprule
Building Scenario & Infiltration Ratio & Indoor PM2.5 (µg/m$^3$) & WHO Exceedance \\
\midrule
Well-sealed (post-2010) & 0.50 & 18.9 & 3.8$\times$ \\
Typical schools (1990--2010) & 0.65 & 24.6 & 4.9$\times$ \\
Poorly-sealed (Soviet-era) & 0.80 & 30.3 & 6.1$\times$ \\
\bottomrule
\end{tabular}
\end{table}

Even well-sealed post-2010 construction (infiltration 0.50) experiences indoor PM2.5 3.8-fold above WHO guidelines. Poorly-sealed Soviet-era buildings---comprising approximately 60\% of Tashkent's school infrastructure---reach 6.1-fold exceedances, creating substantial within-city exposure inequities based on school building age and maintenance~\cite{fisk2017association,allen2016associations}.

Winter indoor exposures showed additional amplification: applying seasonal outdoor means (59.1~µg/m$^3$) to the typical school infiltration ratio (0.65) yields estimated winter classroom concentrations of 38.4~µg/m$^3$---coinciding with peak school occupancy and respiratory infection season (Supplementary Table~S2).

\subsection{Health Impact Assessment Projects Substantial Disease Burden}

Health impact assessment employed a screening-level approach using exposure-response functions from meta-analyses of air pollution epidemiology~\cite{burnett2014integrated,cohen2017estimates}, population data, baseline incidence variation, and exposure measurement. This approach provides conservative, policy-relevant burden estimates summarized in Table~\ref{tab:health}.

\begin{table}[H]
\centering
\caption{Projected annual attributable respiratory risk among school-age children in Tashkent. Estimates derived from screening-level health impact assessment using log-linear exposure-response functions.}
\label{tab:health}
\begin{tabular}{lcc}
\toprule
Parameter & Central Estimate & 95\% CI/Range \\
\midrule
School-age population (5--18 years) & 420,000 & 380,000--460,000 \\
Baseline respiratory incidence & 10\% & 8--15\% \\
Relative risk per 10~µg/m$^3$ & 1.08 & 1.04--1.12 \\
Excess PM2.5 above WHO guideline (5~µg/m$^3$) & 32.9~µg/m$^3$ & --- \\
\addlinespace
\textbf{Relative risk at observed exposure} & \textbf{1.29} & --- \\
\textbf{Attributable fraction (approximate)} & \textbf{22.3\%} & 15--30\% \\
\bottomrule
\end{tabular}
\end{table}

The health impact estimates indicate substantial burden; however, given uncertainties in extrapolating exposure-response functions to Central Asian populations, we present attributable fractions rather than absolute case counts~\cite{anenberg2019global}. Monte Carlo uncertainty propagation (10,000 iterations) sampling from documented parameter uncertainty distributions---relative risk (log-normal, 95\% CI: 1.04--1.12), population (normal, SD: 20,000), baseline incidence (uniform, 8--15\%), and exposure measurement uncertainty (normal, SE: 1.63~µg/m$^3$ representing standard error of the annual mean)---yielded attributable fraction estimates ranging 15--30\% across simulation iterations. Sensitivity analysis varying the counterfactual exposure level (5--15~µg/m$^3$) altered attributable fraction estimates by $\pm$25\%, but critically, intervention ranking and cost-effectiveness conclusions remained unchanged across all parameter scenarios (Supplementary Table~S3). This robustness to parameter uncertainty strengthens confidence in policy recommendations despite limited Central Asian epidemiological data.

The respiratory attributable fraction above uses a deliberately conservative log-linear function (RR~1.08 per 10~µg/m$^3$) appropriate for all-cause respiratory outcomes. Applying paediatric-asthma-specific exposure--response functions yields substantially larger fractions (Fig.~\ref{fig:paferfs}). Using the traffic-related air pollution meta-analysis of Khreis et al.\ (RR~1.03 per 1~µg/m$^3$)~\cite{khreis2017exposure} and the paediatric incidence functions of Anenberg et al.\ (RR~1.34--1.93 per 10~µg/m$^3$)~\cite{anenberg2019global}---and capping the exposure increment at the upper bound of the epidemiological derivation range to limit extrapolation---the attributable fraction for childhood asthma incidence reaches approximately 58\% (95\% CI 26--76\%) and 76\% (59--85\%), respectively. These magnitudes are consistent with global estimates that more than half of paediatric asthma incidence is attributable to anthropogenic PM2.5 in comparably polluted settings such as India and China~\cite{anenberg2019global}. We therefore interpret the 22.3\% respiratory attributable fraction as a conservative floor: the child-specific asthma burden is plausibly two- to three-fold higher, which strengthens rather than weakens the case for protective intervention.

\begin{figure}[H]
\centering
\includegraphics[width=0.7\textwidth]{fig_paf_erfs.png}
\caption{Population attributable fraction of childhood disease from PM2.5 in Tashkent (annual mean 37.9~µg/m$^3$; counterfactual 5~µg/m$^3$), under a conservative all-respiratory function and two paediatric-asthma exposure--response functions (increment capped at the derivation ceiling). Error bars are 95\% intervals from 10$^5$ Monte Carlo draws (seed 42).}
\label{fig:paferfs}
\end{figure}

\subsection{HEPA Filtration Demonstrates Cost-Effective Intervention Potential}

Classroom HEPA filtration intervention modeling (80\% particle removal efficiency~\cite{fisk2017association,qi2017large}) projected substantial exposure reductions. Cost-effectiveness analysis for citywide deployment assumed: 800 Tashkent schools, 25 classrooms per school, \$2,500 capital cost per classroom unit, 5-year equipment lifespan, and 10\% annual maintenance costs (Table~\ref{tab:cost}).

\begin{table}[H]
\centering
\caption{Cost-effectiveness of citywide classroom HEPA filtration deployment in Tashkent schools.}
\label{tab:cost}
\begin{tabular}{lc}
\toprule
Parameter & Value \\
\midrule
Schools in Tashkent & 800 \\
Classrooms per school (average) & 25 \\
Total classrooms & 20,000 \\
Capital cost per classroom & \$2,500 \\
Total capital investment & \$50 million \\
Equipment lifespan & 5 years \\
Annualized capital cost & \$10 million/year \\
Maintenance (10\% annually) & \$5 million/year \\
Total annual cost & \$15 million/year \\
\addlinespace
Projected indoor PM2.5 reduction & 50--80\% \\
Post-intervention indoor exposure & 3.8--15.2~µg/m$^3$ \\
\bottomrule
\end{tabular}
\end{table}

At these intervention costs, classroom HEPA filtration compares favorably with established global health interventions~\cite{jamison2017disease,allen2016associations}. A detailed sensitivity analysis of cost-effectiveness parameters is provided in Supplementary Table~S5. From an urban sustainability and equity perspective, this framework directly addresses Sustainable Development Goal 11 (safe, inclusive cities and communities), SDG 4 (quality education and learning environments), and SDG 3 (health and wellbeing), demonstrating how school-based interventions can serve as co-benefit strategies linking urban planning, education, and health systems in rapidly urbanizing regions.

\subsection{Regional Generalization: A Six-Capital Comparison}

To test whether the framework generalizes and to place Tashkent in regional context, we applied the identical pipeline to the U.S. Embassy reference monitors in all six Central Asian capitals over the same window (Table~\ref{tab:regional}, Fig.~\ref{fig:multicity}). Five of the six capitals exceeded the WHO annual guideline 4.6--10.7-fold. Dushanbe was the most polluted (annual 53.3~µg/m$^3$; school-hours 52.1; 97\% of days above the 24-hour guideline), followed by Tashkent (37.9), Bishkek (35.6), and Almaty (34.2); Ashgabat (22.8) and Astana (18.5; partial 2022--2023 coverage) were lower but still well above guideline. School-hours concentrations tracked annual means closely (17.5--52.1~µg/m$^3$), and every city showed pronounced winter amplification driven by heating and inversions (winter means 28--84~µg/m$^3$). Applying the conservative respiratory function, the population attributable fraction for school-age children ranged from 10\% (Astana) to 31\% (Dushanbe). The consistency of these patterns---high winter-dominated exposure during peak school occupancy across five capitals---demonstrates that the school-exposure problem is regional rather than Tashkent-specific, and that a single open reference monitor per city is sufficient to flag it. This is the core transferability claim of the framework: the same minimal inputs that characterized Tashkent immediately extend to its neighbours, identifying Dushanbe in particular as warranting priority attention. Extending the full spatial layer to Dushanbe confirms the pattern: across 178 mapped schools, modelled indoor classroom concentrations span 26.7--46.1~µg/m$^3$ (Supplementary Fig.~S7), so that even the best-sealed Dushanbe classroom exceeds the WHO annual guideline more than fivefold and the worst by ninefold---a larger absolute burden than any Tashkent school.

\begin{table}[H]
\centering
\caption{School-age PM2.5 exposure across six Central Asian capitals from U.S. Embassy reference monitors (OpenAQ; January 2022--June 2023). PAF is the conservative all-respiratory population attributable fraction (RR 1.08 per 10~µg/m$^3$, counterfactual 5~µg/m$^3$); typical-classroom is the annual mean times a 0.65 infiltration factor. $^{\dagger}$Astana had partial coverage (77 valid days) and is reported for completeness only.}
\label{tab:regional}
\small
\begin{tabular}{lcccccc}
\toprule
City & Annual & School-hrs & Winter & Days $>$15 & PAF & Typical \\
 & (µg/m$^3$) & (µg/m$^3$) & (µg/m$^3$) & (\%) & (\%) & classroom \\
\midrule
Dushanbe & 53.3 & 52.1 & 84.2 & 96.7 & 31.0 & 34.6 \\
\textbf{Tashkent} & \textbf{37.9} & \textbf{33.7} & \textbf{59.1} & \textbf{93.1} & \textbf{22.4} & \textbf{24.6} \\
Bishkek & 35.6 & 35.7 & 79.2 & 68.1 & 21.0 & 23.1 \\
Almaty & 34.2 & 39.5 & 66.7 & 70.1 & 20.2 & 22.3 \\
Ashgabat & 22.8 & 23.1 & 28.3 & 66.1 & 12.8 & 14.8 \\
Astana$^{\dagger}$ & 18.5 & 17.5 & 31.3 & 41.6 & 9.9 & 12.1 \\
\bottomrule
\end{tabular}
\end{table}

\begin{figure}[H]
\centering
\includegraphics[width=0.85\textwidth]{fig_multicity.png}
\caption{Annual-mean and school-hours PM2.5 across six Central Asian capitals (U.S. Embassy reference monitors, January 2022--June 2023), with WHO guidelines and the conservative respiratory attributable fraction (PAF) per city. $^{*}$Astana has partial coverage (77 valid days).}
\label{fig:multicity}
\end{figure}

%----------------------------------------------------------------------
\section{Discussion}
%----------------------------------------------------------------------

This study reframes Tashkent's well-documented air-pollution crisis around the population most exposed during its developmental years---school-age children---and the building environments in which they spend their days. Rather than re-deriving the city's ambient burden, which the World Bank's 2024 chemical-transport assessment has authoritatively quantified (population-weighted annual mean 38.8~µg/m$^3$; approximately 3,000 attributable premature deaths; US\$488~million in annual welfare loss)~\cite{worldbank2024tashkent}, we add the exposure-and-protection layer that ambient assessments omit: where schools sit, how classroom envelopes modify exposure, what share of childhood disease is attributable, and which near-term protections are justified. Our screening-level burden estimate is consistent with the World Bank's independent figure, which we adopt as the reference baseline; we do not treat the numerical proximity as validation, since both ultimately rest on the same reference-monitor record.

The exposure documented here is severe and chronic rather than episodic. An annual mean of 37.9~µg/m$^3$ (7.6 times the WHO guideline), 93.1\% of days above the 24-hour guideline, and a 126\% winter amplification place Tashkent above Beijing's post-intervention concentrations (29--33~µg/m$^3$)~\cite{crea2023china}. The temporal structure points away from local traffic: weekday and weekend school-hour concentrations are indistinguishable, overnight levels rival the morning peak, and the World Bank source apportionment attributes only 16\% of ambient concentration to transport against 28\% to residential heating and 36\% to windblown dust~\cite{worldbank2024tashkent}. Because the dominant sources are diffuse, seasonal, and partly transboundary, they cannot be controlled on the timescale relevant to children currently in school---which shifts the locus of near-term protection from the ambient airshed to the classroom.

Two within-city dimensions determine how that ambient burden reaches children. First, although the basin produces a comparatively uniform regional field, schools are not uniformly placed: 15.4\% of the 605 schools we mapped lie within 100~m of a major road and a quarter within 200~m, where traffic and resuspended-dust increments concentrate~\cite{apte2017high,khreis2017exposure}. Second, and more consequentially, the building envelope governs how much outdoor PM2.5 penetrates the classroom. Across the school stock, modelled indoor concentrations span 18.9--30.3~µg/m$^3$---a 1.6-fold gradient---so that children in poorly sealed Soviet-era buildings (about 60\% of the stock) experience indoor exposures 60\% higher than peers in post-2010 construction, an equity gradient aligned with building age and, by extension, neighbourhood resourcing. Even the best-sealed classrooms remain at 3.8 times the WHO guideline.

The health stakes for children are correspondingly large. Our conservative respiratory attributable fraction (22.3\%) rests on an all-cause respiratory slope; applying paediatric-asthma-specific functions raises the attributable fraction for childhood asthma incidence to roughly 58--76\%, consistent with global estimates that more than half of paediatric asthma in comparably polluted countries is attributable to anthropogenic PM2.5~\cite{khreis2017exposure,anenberg2019global}. We present these as a graded set---a conservative floor and a child-specific ceiling---rather than a single point estimate; the policy implication is robust to that uncertainty, since the burden borne by children is substantial under any defensible function.

Framed this way, classroom filtration is not merely a health measure but a human-capital one. Observational studies already link traffic-related pollution at schools to slower cognitive development in children~\cite{sunyer2015association}, and randomized trials now connect reductions in classroom PM2.5 to measurable cognitive and educational gains: a trial in Lahore found a 0.15 standard-deviation improvement in test scores, concentrated in mathematics and fluid reasoning, within weeks of installing purifiers~\cite{malik2024classroom}, and a double-blind crossover in Tianjin found higher standardized-test scores under active filtration~\cite{marshall2024testscores}. Filtration also reduces in-school respiratory infection (relative risk approximately 0.73 in a controlled study)~\cite{banholzer2024aircleaners}, and portable HEPA units typically lower indoor PM2.5 by 40--60\% in classrooms~\cite{carmona2022schools,chen2015effectiveness,allen2011}. The marginal benefit is concentration-dependent: a large Los Angeles trial found little incremental health effect because baseline classroom concentrations were already near 1.5~µg/m$^3$~\cite{vieira2025hifive}, whereas Tashkent classrooms sit an order of magnitude higher---precisely the regime in which the Lahore and Tianjin benefits were observed. At an estimated US\$15~million per year for citywide deployment---about 3\% of the documented annual health cost---filtration is favourable even before its learning co-benefits are counted, and it can be targeted first to the high-infiltration, near-road schools our mapping identifies.

These findings should be read with their limitations. The temporal characterization rests on a single, centrally located reference monitor; while the basin's homogeneity (World Bank inter-station correlation $r>0.85$) supports its use for citywide temporal burden, it cannot resolve near-road or indoor microenvironments, which warrant direct school monitoring~\cite{wilson2005spatial,jerrett2005spatial}. The 18-month record spans multiple heating seasons but not full inter-annual variability. Classroom concentrations are modelled from literature infiltration factors rather than measured in situ~\cite{abt2000relative,chen2012review,macintosh2009residential}, and school locations are drawn from OpenStreetMap, which is near-complete for Tashkent but not authoritative. Our exposure-response functions are transferred from North American and European cohorts and, for the higher childhood-asthma estimates, extrapolated toward the upper bound of their derivation range; we therefore report attributable fractions rather than case counts~\cite{burnett2014integrated,cohen2017estimates}. Finally, an earlier phase of this work used a low-cost optical sensor that over-read the reference monitor by roughly 50\% (annual mean 56 versus 38~µg/m$^3$) under high-humidity winter conditions; propagated through the same burden model, this single data-source choice would raise the respiratory attributable fraction from 22\% to 33\% and inflate estimated health costs commensurately. That experience is itself a governance lesson: in data-scarce cities increasingly reliant on low-cost or citizen networks, anchoring exposure assessment to a reference-grade instrument is essential before such estimates inform policy~\cite{morawska2019applications}.

The broader contribution is a transferable template. Combining one open reference monitor, open building-stock and school-location data, and published exposure-response and infiltration functions, the workflow converts minimal inputs into a spatially explicit, equity-aware, child-focused exposure and intervention assessment, without new monitoring infrastructure. U.S. diplomatic missions operate reference monitors in most capitals, and the same open data layers exist globally, so the approach extends directly to other Central Asian capitals---Almaty, Bishkek, Dushanbe, Ashgabat---and to data-scarce cities more widely~\cite{worldbank2024centralasia,tursumbayeva2023}. For Tashkent specifically, the analysis supports four actions: prioritising classroom filtration in high-infiltration, near-road schools; incorporating major-road proximity into school-siting and renovation standards; piloting arrival- and dismissal-time traffic management at the near-road schools identified here; and extending direct monitoring into classrooms to convert this baseline into routine, equity-targeted management. Connecting air quality, education, and equity, these measures advance Sustainable Development Goals 3, 4, and 11 together~\cite{nieuwenhuijsen2019influence,landrigan2017lancet}.

%----------------------------------------------------------------------
\section{Methods}
%----------------------------------------------------------------------

\subsection{Study Design and Setting}

We conducted an observational ecological time-series study examining PM2.5 exposure and projected health outcomes among school-age children in Tashkent, Uzbekistan. Tashkent (41.32°N, 69.27°E; metropolitan population 2.6 million) is located in an enclosed basin characterized by continental climate, winter thermal inversions (mixing layer 300--400~m), and limited precipitation (400~mm annually). The city hosts approximately 800 schools serving 420,000 children aged 5--18 years.

\subsection{Data Source and Quality Assurance}

PM2.5 measurements were obtained from the OpenAQ platform aggregating government air quality monitoring data. Station 8881 (U.S. Embassy Tashkent), operated by the U.S. Department of State under the StateAir program, uses a Federal Equivalent Method (FEM) beta attenuation monitor with hourly temporal resolution. This monitoring station provides research-grade measurements calibrated to U.S. EPA standards and is considered the reference monitor for Tashkent air quality. The 18-month study period (January 1, 2022--June 29, 2023; 545 days) yielded 8,301 valid hourly measurements from a possible 13,080 hours (63.5\% coverage). For transparency, we note that an exploratory phase of this project evaluated a co-located low-cost optical sensor; all analyses reported here use exclusively the reference-grade FEM record from Station 8881, and the optical comparison is retained only as a data-governance caveat (Discussion). For the six-capital comparison (Table~\ref{tab:regional}), we applied the identical retrieval and quality-control pipeline to the U.S. Embassy reference monitors in Almaty (OpenAQ location 8876), Astana (7094), Bishkek (8827), Dushanbe (8684), and Ashgabat (8870) over the same period, using the same $\geq$18-hour daily-completeness threshold; Astana returned only 77 valid days and is reported for completeness only.

Quality control procedures employed a two-tier validation system: (1) \textit{instrument validity bounds} (0--1000~µg/m$^3$) reflecting the BAM-1020 measurement range, with values outside these bounds excluded as sensor errors; (2) \textit{extreme-event flagging} ($>$400~µg/m$^3$) requiring manual review---12 flagged records (0.02\%) were retained as valid winter pollution episodes consistent with documented Central Asian thermal inversions. Additional QC included exclusion of null/missing values and negative readings (5.8\% of records), outlier review using IQR method (1.5$\times$IQR threshold) with retention after verification of extreme winter values, and daily aggregation requiring $\geq$18 valid hourly observations (75\% completeness threshold). Missing observations occurred primarily during sensor maintenance periods, with no systematic temporal bias detected (Supplementary Figure~S4). Measurement uncertainty sources and sensor characteristics are documented in Supplementary Table~S6.

\subsection{Exposure Assessment}

Primary exposure metric was 24-hour mean PM2.5 concentration (µg/m$^3$). Measurements were classified against WHO 2021 Air Quality Guidelines: annual guideline 5~µg/m$^3$; 24-hour guideline 15~µg/m$^3$ (99th percentile)~\cite{who2021guidelines}.

Indoor school exposure was modeled using infiltration factors (indoor/outdoor ratios 0.50--0.80) derived from published Central Asian building stock literature~\cite{abt2000relative,macintosh2009residential,chen2012review}. School exposure windows were defined as 08:00--15:00 weekdays; commute windows as 07:00--09:00 and 16:00--18:00.

Seasonal classification followed calendar months: winter (December--February), spring (March--May), summer (June--August), fall (September--November).

\subsection{Spatial School Analysis}

School locations were obtained from OpenStreetMap (OSM; \texttt{amenity=school}) within the Tashkent metropolitan bounding box (41.15--41.42$^{\circ}$N, 69.10--69.45$^{\circ}$E), yielding 605 geolocated schools. Major roads (OSM \texttt{highway} classes motorway, trunk, and primary) were retrieved for the same extent. Coordinates were projected to a local metric frame, and each school's Euclidean distance to the nearest major-road segment was computed using an R-tree spatial index; schools within 100~m and 200~m were flagged as near-road microenvironments. The regional background was anchored to the reference-monitor annual mean and checked for spatial homogeneity against satellite-derived ACAG estimates and the World Bank CAMx inter-station correlation ($r>0.85$); a 25-point grid of CAMS-global reanalysis annual means (retrieved via the Open-Meteo Air Quality API for 2022--2023) confirmed that coarse global products do not resolve intra-urban gradients here (two distinct cells across the metro; $\sim$50\% under-read of the reference monitor). The outdoor exposure surface added a near-road increment, $\mathrm{inc}(d)=A\,e^{-d/L}$ ($A=5$~µg/m$^3$, $L=150$~m; sensitivity $A\in\{3,5,8\}$, $L\in\{100,200\}$~m), to this background. Building-envelope heterogeneity was represented by assigning each school a construction-era infiltration factor sampled from the documented stock composition (60\% Soviet-era, 0.80; 30\% 1990--2010, 0.65; 10\% post-2010, 0.50); indoor classroom concentrations were obtained by multiplying the outdoor surface by this factor. These are first-order proximity-model estimates pending in-classroom measurement; all spatial code and retrieved grids are archived with the analysis scripts.

\subsection{Health Impact Assessment}

Attributable respiratory burden was estimated using the population attributable fraction approach with Global Burden of Disease Integrated Exposure-Response (IER) functions and the associated global exposure framework~\cite{burnett2014integrated,cohen2017estimates,brauer2016ambient,shaddick2018data,lelieveld2015contribution}:

\begin{equation}
\text{AF} = \frac{RR^{(C-C_0)/10} - 1}{RR^{(C-C_0)/10}}
\end{equation}

\noindent where $C$ = observed PM2.5 concentration, $C_0$ = counterfactual level (5~µg/m$^3$, WHO annual guideline), and $RR$ = relative risk per 10~µg/m$^3$ increment (1.08; 95\% CI: 1.04--1.12) for pediatric respiratory outcomes.

Monte Carlo uncertainty propagation (10,000 iterations; random seed=42) sampled from: RR (log-normal distribution based on 95\% CI); population (normal; mean 420,000, SD 20,000); baseline incidence (uniform; 8--15\%); exposure (normal; mean 37.9, SE 1.63~µg/m$^3$ representing standard error of the annual mean---distinct from SD 38.4~µg/m$^3$ describing hourly concentration variability).

\textbf{Uncertainty decomposition:} The SE of 1.63~µg/m$^3$ reflects sampling uncertainty on the annual mean (i.e., precision of the central estimate), while the SD of 38.4~µg/m$^3$ reflects day-to-day PM2.5 variability driven by meteorology and emission patterns. Additionally, systematic measurement uncertainty from the BAM-1020 instrument (±7.8~µg/m$^3$ combined from calibration, temperature, and humidity effects; Supplementary Table~S6) was incorporated through sensitivity analysis rather than the primary Monte Carlo simulation, as this uncertainty affects the mean estimate location rather than parameter sampling distributions.

\textbf{Paediatric asthma functions.} In addition to the conservative all-respiratory function (RR 1.08 per 10~µg/m$^3$), we computed attributable fractions for childhood asthma incidence using two paediatric exposure--response functions: Khreis et al.\ (2017), RR 1.03 per 1~µg/m$^3$ (95\% CI 1.01--1.05), and Anenberg et al.\ (2018), RR 1.34--1.93 per 10~µg/m$^3$~\cite{khreis2017exposure,anenberg2019global}. To limit extrapolation beyond the epidemiological derivation range, the exposure increment was capped at 29~µg/m$^3$ (observed concentration capped at $\sim$34~µg/m$^3$) for these functions. Attributable fractions and 95\% intervals were obtained from $10^5$ Monte Carlo draws (seed 42).

\subsection{Intervention Cost-Effectiveness}

HEPA filtration cost-effectiveness modeling assumed: 800 Tashkent schools; 25 classrooms/school; \$2,500 capital cost/classroom; 80\% particle removal efficiency~\cite{fisk2017association,qi2017large}; 5-year equipment life; 10\% annual maintenance. Annualized costs were compared with projected cases prevented. DALY estimates assumed 0.15 disability weight for respiratory symptoms over 7-day average duration.

\subsection{Statistical Analysis}

Analyses used Python 3.12 with pandas 2.0, numpy 1.24, scipy 1.10, and matplotlib 3.7. Descriptive statistics included mean, median, SD, IQR, and percentiles. Seasonal comparisons used Kruskal-Wallis H-test ($\alpha$=0.05). Uncertainty intervals are 95\% credible intervals from Monte Carlo simulation unless otherwise specified.

\subsection{Ethics Statement}

This study used exclusively public environmental monitoring data from OpenAQ containing no human participants, personal information, or individual health records. The study qualifies for institutional review board exemption under 45 CFR 46.104(d) as research involving existing public data without personal identifiers.

%----------------------------------------------------------------------
\section*{Data Availability}
%----------------------------------------------------------------------

All data supporting this study are publicly available. Raw PM2.5 measurements: OpenAQ API v3, from the U.S. Embassy reference monitors in Tashkent (location 8881), Almaty (8876), Astana (7094), Bishkek (8827), Dushanbe (8684), and Ashgabat (8870) (e.g.\ https://api.openaq.org/v3/locations/8881). Complete reproducibility dataset including processed data, analysis scripts, and output files: Zenodo concept DOI 10.5281/zenodo.17792118 (resolving to the current reference-grade version) and GitHub repository https://github.com/ProgrmerJack/Air-quality-insight-Uzbekistan (MIT License). School and road geodata: OpenStreetMap. Satellite PM2.5 cross-check: Atmospheric Composition Analysis Group (ACAG). WHO 2021 Air Quality Guidelines: https://www.who.int/publications/i/item/9789240034228. GBD 2021 exposure-response functions: Burnett RT, Pope CA, Ezzati M, et al. An integrated risk function for estimating the global burden of disease attributable to ambient fine particulate matter exposure. Environ Health Perspect. 2014;122(4):397-403. doi:10.1289/ehp.1307049. U.S. Department of State Air Quality Monitoring Program (StateAir): https://www.airnow.gov/international/us-embassies-and-consulates/.

%----------------------------------------------------------------------
\section*{Code Availability}
%----------------------------------------------------------------------

Complete analysis code (Python) including data processing, statistical analysis, indoor infiltration modeling, and visualization is available at https://github.com/ProgrmerJack/Air-quality-insight-Uzbekistan under MIT License.

%----------------------------------------------------------------------
\section*{Author Contributions}
%----------------------------------------------------------------------

A.A. conceived the study, obtained and processed data, performed all analyses, created visualizations, and wrote the manuscript.

%----------------------------------------------------------------------
\section*{Competing Interests}
%----------------------------------------------------------------------

The author declares no competing financial or non-financial interests.

%----------------------------------------------------------------------
\section*{Acknowledgments}
%----------------------------------------------------------------------

This research received no specific grant from any funding agency in the public, commercial, or not-for-profit sectors. The author acknowledges OpenAQ for providing open access to real-time air quality data, the U.S. Department of State for operating the StateAir monitoring program that provided reference-grade PM2.5 measurements, OpenStreetMap contributors for the school and road geodata, and the Ministry of Ecology, Environmental Protection and Climate Change of the Republic of Uzbekistan for supporting environmental monitoring initiatives.

%----------------------------------------------------------------------
% REFERENCES - Expanded to 36 citations including World Bank/UNECE policy documents
%----------------------------------------------------------------------

\bibliographystyle{unsrtnat}
\begin{thebibliography}{36}

\bibitem[Abt et al.(2000)]{abt2000relative}
Abt, E., Suh, H. H., Allen, G., \& Koutrakis, P. (2000). Relative contribution of outdoor and indoor particle sources to indoor concentrations. \textit{Environmental Science \& Technology}, 34(17), 3579--3587.

\bibitem[Allen et al.(2011)]{allen2011}
Allen, R. W., Carlsten, C., Karlen, B., Leckie, S., van Eeden, S., Carlsten, J., \& Brauer, M. (2011). An air filter intervention study of endothelial function among healthy adults in a woodsmoke-impacted community. \textit{American Journal of Respiratory and Critical Care Medicine}, 183(9), 1378--1386.

\bibitem[Allen et al.(2016)]{allen2016associations}
Allen, J. G., MacNaughton, P., Satish, U., Santanam, S., Vallarino, J., \& Spengler, J. D. (2016). Associations of cognitive function scores with carbon dioxide, ventilation, and volatile organic compound exposures in office workers: A controlled exposure study of green and conventional office environments. \textit{Environmental Health Perspectives}, 124(6), 805--812.

\bibitem[Anenberg et al.(2018)]{anenberg2019global}
Anenberg, S. C., Henze, D. K., Tinney, V., Kinney, P. L., Raich, W., Fann, N., Malley, C. S., Roman, H., Lamsal, L., Duncan, B., Martin, R. V., van Donkelaar, A., Brauer, M., Doherty, R., Jonson, J. E., Davila, Y., Sudo, K., \& Kuylenstierna, J. C. I. (2018). Estimates of the global burden of ambient PM2.5, ozone, and NO2 on asthma incidence and emergency room visits. \textit{Environmental Health Perspectives}, 126(10), 107004.

\bibitem[Apte et al.(2017)]{apte2017high}
Apte, J. S., Messier, K. P., Gani, S., Brauer, M., Kirchstetter, T. W., Lunden, M. M., Marshall, J. D., Portier, C. J., Vermeulen, R. C. H., \& Hamburg, S. P. (2017). High-resolution air pollution mapping with Google Street View cars. \textit{Environmental Science \& Technology}, 51(12), 6999--7008.

\bibitem[Ashuraliyev et al.(2025)]{ashuraliyev2025pm25tashkent}
Ashuraliyev, A., Qudayberganov, T., \& Mamatov, F. (2025). Particulate air pollution in Tashkent and consequences to population health. \textit{Russian Journal of Social Medicine}, 104, 1.

\bibitem[Brauer et al.(2016)]{brauer2016ambient}
Brauer, M., Freedman, G., Frostad, J., van Donkelaar, A., Martin, R. V., Dentener, F., van Dingenen, R., Estep, K., Amini, H., Apte, J. S., Balakrishnan, K., Barregard, L., Broday, D., Feigin, V., Ghosh, S., Hopke, P. K., Knibbs, L. D., Kokubo, Y., Liu, Y., ... Cohen, A. (2016). Ambient air pollution exposure estimation for the Global Burden of Disease 2013. \textit{Environmental Science \& Technology}, 50(1), 79--88.

\bibitem[Brook et al.(2010)]{brook2010particulate}
Brook, R. D., Rajagopalan, S., Pope, C. A., Brook, J. R., Bhatnagar, A., Diez-Roux, A. V., Holguin, F., Hong, Y., Luepker, R. V., Mittleman, M. A., Peters, A., Siscovick, D., Smith, S. C., Whitsel, L., \& Kaufman, J. D. (2010). Particulate matter air pollution and cardiovascular disease: An update to the scientific statement from the American Heart Association. \textit{Circulation}, 121(21), 2331--2378.

\bibitem[Burnett et al.(2014)]{burnett2014integrated}
Burnett, R. T., Pope, C. A., Ezzati, M., Olives, C., Lim, S. S., Mehta, S., Shin, H. H., Singh, G., Hubbell, B., Brauer, M., Anderson, H. R., Smith, K. R., Balmes, J. R., Bruce, N. G., Kan, H., Laden, F., Pr\"uss-Ust\"un, A., Turner, M. C., Gapstur, S. M., ... Cohen, A. (2014). An integrated risk function for estimating the global burden of disease attributable to ambient fine particulate matter exposure. \textit{Environmental Health Perspectives}, 122(4), 397--403.

\bibitem[Carmona et al.(2022)]{carmona2022schools}
Carmona, N., Seto, E., Gould, T. R., Rathbun, S. L., Austin, E., Shirai, J. H., Spalt, E. W., \& Larson, T. V. (2022). Indoor air quality intervention in schools: Effectiveness of a portable HEPA filter deployment in five schools impacted by roadway and aircraft pollution sources. \textit{Atmosphere}, 13(10), 1623.

\bibitem[Chen \& Zhao(2012)]{chen2012review}
Chen, C., \& Zhao, B. (2012). Review of relationship between indoor and outdoor particles: I/O ratio, infiltration factor and penetration factor. \textit{Atmospheric Environment}, 49, 175--189.

\bibitem[Chen et al.(2015)]{chen2015effectiveness}
Chen, R., Zhao, A., Chen, H., Zhao, Z., Cai, J., Wang, C., Yang, C., Li, H., Xu, X., Ha, S., Li, T., \& Kan, H. (2015). Cardiopulmonary benefits of reducing indoor particles of outdoor origin: A randomized, double-blind crossover trial of air purifiers. \textit{Journal of the American College of Cardiology}, 65(21), 2279--2287.

\bibitem[Cohen et al.(2017)]{cohen2017estimates}
Cohen, A. J., Brauer, M., Burnett, R., Anderson, H. R., Frostad, J., Estep, K., Balakrishnan, K., Brunekreef, B., Dandona, L., Dandona, R., Feigin, V., Freedman, G., Hubbell, B., Jobling, A., Kan, H., Knibbs, L., Liu, Y., Martin, R., Morawska, L., ... Forouzanfar, M. H. (2017). Estimates and 25-year trends of the global burden of disease attributable to ambient air pollution. \textit{The Lancet}, 389(10082), 1907--1918.

\bibitem[CREA(2023)]{crea2023china}
Centre for Research on Energy and Clean Air. (2023). PM2.5 rebounds in China in 2023, after falling for 10 years straight. CREA Analysis.

\bibitem[Fisk(2013)]{fisk2017association}
Fisk, W. J. (2013). Health benefits of particle filtration. \textit{Indoor Air}, 23(5), 357--368.

\bibitem[Gauderman et al.(2015)]{gauderman2015association}
Gauderman, W. J., Urman, R., Avol, E., Berhane, K., McConnell, R., Rappaport, E., Chang, R., Lurmann, F., \& Gilliland, F. (2015). Association of improved air quality with lung development in children. \textit{New England Journal of Medicine}, 372(19), 1803--1813.

\bibitem[Gehring et al.(2013)]{gehring2013air}
Gehring, U., Gruzieva, O., Agius, R. M., Beelen, R., Custovic, A., Cyrys, J., Eeftens, M., Flexeder, C., Fuertes, E., Heinrich, J., Hoffmann, B., de Jongste, J. C., Kerkhof, M., Klumper, C., Korek, M., Molter, A., Schultz, E. S., Simpson, A., Sugiri, D., ... Brunekreef, B. (2013). Air pollution exposure and lung function in children: The ESCAPE project. \textit{Environmental Health Perspectives}, 121(11-12), 1357--1364.

\bibitem[Hammer et al.(2020)]{acag2022satpm25}
Hammer, M. S., van Donkelaar, A., Li, C., Lyapustin, A., Sayer, A. M., Hsu, N. C., Levy, R. C., Garay, M. J., Kalashnikova, O. V., Kahn, R. A., Brauer, M., Apte, J. S., Henze, D. K., Zhang, L., Zhang, Q., Ford, B., Pierce, J. R., \& Martin, R. V. (2020). Global estimates and long-term trends of fine particulate matter concentrations (1998--2018). \textit{Environmental Science \& Technology}, 54(13), 7879--7890.

\bibitem[Jamison et al.(2017)]{jamison2017disease}
Jamison, D. T., Gelband, H., Horton, S., Jha, P., Laxminarayan, R., Mock, C. N., \& Nugent, R. (Eds.). (2017). \textit{Disease Control Priorities: Improving Health and Reducing Poverty} (3rd ed.). World Bank.

\bibitem[Jerrett et al.(2005)]{jerrett2005spatial}
Jerrett, M., Burnett, R. T., Ma, R., Pope, C. A., Krewski, D., Newbold, K. B., Thurston, G., Shi, Y., Finkelstein, N., Calle, E. E., \& Thun, M. J. (2005). Spatial analysis of air pollution and mortality in Los Angeles. \textit{Epidemiology}, 16(6), 727--736.

\bibitem[Khreis et al.(2017)]{khreis2017exposure}
Khreis, H., Kelly, C., Tate, J., Parslow, R., Lucas, K., \& Nieuwenhuijsen, M. (2017). Exposure to traffic-related air pollution and risk of development of childhood asthma: A systematic review and meta-analysis. \textit{Environment International}, 100, 1--31.

\bibitem[Landrigan et al.(2018)]{landrigan2017lancet}
Landrigan, P. J., Fuller, R., Acosta, N. J. R., Adeyi, O., Arnold, R., Basu, N., Bald\'e, A. B., Bertollini, R., Bose-O'Reilly, S., Boufford, J. I., Breysse, P. N., Chiles, T., Mahidol, C., Coll-Seck, A. M., Cropper, M. L., Fobil, J., Fuster, V., Greenstone, M., Haines, A., ... Zhong, M. (2018). The Lancet Commission on pollution and health. \textit{The Lancet}, 391(10119), 462--512.

\bibitem[Lelieveld et al.(2015)]{lelieveld2015contribution}
Lelieveld, J., Evans, J. S., Fnais, M., Giannadaki, D., \& Pozzer, A. (2015). The contribution of outdoor air pollution sources to premature mortality on a global scale. \textit{Nature}, 525(7569), 367--371.

\bibitem[Li et al.(2019)]{li2017wintertime}
Li, K., Jacob, D. J., Liao, H., Shen, L., Zhang, Q., \& Bates, K. H. (2019). A two-pollutant strategy for improving ozone and particulate air quality in China. \textit{Nature Geoscience}, 12(11), 906--910.

\bibitem[MacIntosh et al.(2009)]{macintosh2009residential}
MacIntosh, D. L., Minegishi, T., Kaufman, M., Baker, B. J., Allen, J. G., Levy, J. I., \& Myatt, T. A. (2009). The benefits of whole-house in-duct air cleaning in reducing exposures to fine particulate matter of outdoor origin. \textit{Journal of Exposure Science \& Environmental Epidemiology}, 19(6), 628--637.

\bibitem[Morawska et al.(2018)]{morawska2019applications}
Morawska, L., Thai, P. K., Liu, X., Asumadu-Sakyi, A., Ayoko, G., Bartonova, A., Bedber, A., Chai, F., Christensen, B., Dunbabin, M., Gao, J., Hagler, G. S. W., Jayaratne, R., Kumar, P., Lau, A. K. H., Louie, P. K. K., Mazaheri, M., Ning, Z., Motta, N., ... Williams, R. (2018). Applications of low-cost sensing technologies for air quality monitoring and exposure assessment. \textit{Environment International}, 116, 286--299.

\bibitem[Nieuwenhuijsen(2018)]{nieuwenhuijsen2019influence}
Nieuwenhuijsen, M. J. (2018). Influence of urban and transport planning and the city environment on cardiovascular disease. \textit{Nature Reviews Cardiology}, 15(7), 432--438.

\bibitem[Pope et al.(2004)]{pope2004cardiovascular}
Pope, C. A., Burnett, R. T., Thurston, G. D., Thun, M. J., Calle, E. E., Krewski, D., \& Godleski, J. J. (2004). Cardiovascular mortality and long-term exposure to particulate air pollution. \textit{Circulation}, 109(1), 71--77.

\bibitem[Qi et al.(2017)]{qi2017large}
Qi, Y., Qian, H., Sun, L., Li, P., Li, J., Ma, L., Li, H., Li, Y., \& Li, Q. (2017). Impacts of air cleaners on cognitive function in patients with mild-to-moderate Alzheimer's disease: A pilot intervention study. \textit{Journal of Alzheimer's Disease}, 59(3), 1055--1067.

\bibitem[Schwartz(2004)]{schwartz2004air}
Schwartz, J. (2004). Air pollution and children's health. \textit{Pediatrics}, 113(S3), 1037--1043.

\bibitem[Shaddick et al.(2018)]{shaddick2018data}
Shaddick, G., Thomas, M. L., Amini, H., Brober, D., Cohen, A., Frostad, J., Green, A., Gumy, S., Liu, Y., Martin, R. V., Pruss-Ustun, A., Simpson, D., van Donkelaar, A., \& Brauer, M. (2018). Data integration for the assessment of population exposure to ambient air pollution for global burden of disease assessment. \textit{Environmental Science \& Technology}, 52(16), 9069--9078.

\bibitem[Sunyer et al.(2015)]{sunyer2015association}
Sunyer, J., Esnaola, M., Alvarez-Pedrerol, M., Forns, J., Rivas, I., L\'opez-Vicente, M., Suades-Gonz\'alez, E., Foraster, M., Garcia-Esteban, R., Basaga\~na, X., Viana, M., Cirach, M., Moreno, T., Alastuey, A., Sebastian-Galles, N., Nieuwenhuijsen, M., \& Querol, X. (2015). Association between traffic-related air pollution in schools and cognitive development in primary school children. \textit{PLoS Medicine}, 12(3), e1001792.

\bibitem[Tursumbayeva et al.(2023)]{tursumbayeva2023}
Tursumbayeva, M., Kalisa, E., Crape, B., \& Kerimray, A. (2023). Cities of Central Asia: New hotspots of air pollution in the world. \textit{Atmospheric Environment}, 305, 119901.

\bibitem[UNECE(2024)]{unece2024uzbekistan}
United Nations Economic Commission for Europe. (2024). \textit{Environmental Performance Reviews: Uzbekistan---Fourth Review}. UNECE.

\bibitem[WHO(2021)]{who2021guidelines}
World Health Organization. (2021). \textit{WHO Global Air Quality Guidelines: Particulate Matter (PM2.5 and PM10), Ozone, Nitrogen Dioxide, Sulfur Dioxide and Carbon Monoxide}. WHO.

\bibitem[Wilson et al.(2005)]{wilson2005spatial}
Wilson, J. G., Kingham, S., Pearce, J., \& Sturman, A. P. (2005). A review of intraurban variations in particulate air pollution: Implications for epidemiological research. \textit{Atmospheric Environment}, 39(34), 6444--6462.

\bibitem[World Bank(2024a)]{worldbank2024centralasia}
World Bank. (2024a). \textit{Air Quality Management in Central Asia: Priorities, Needs, and Solutions}. World Bank.

\bibitem[World Bank(2024b)]{worldbank2024uzbekistan}
World Bank. (2024b). \textit{Uzbekistan Climate Adaptation and Resilience Assessment}. World Bank.

\bibitem[World Bank \& UNEP(2024)]{worldbank2024dialogue}
World Bank \& UNEP. (2024). First High-Level Policy Dialogue on Air Quality Management in Central Asia. Tashkent, Uzbekistan.

\bibitem[Zhang et al.(2019)]{zhang2017drivers}
Zhang, Q., Zheng, Y., Tong, D., Shao, M., Wang, S., Zhang, Y., Xu, X., Wang, J., He, H., Liu, W., Ding, Y., Lei, Y., Li, J., Wang, Z., Zhang, X., Wang, Y., Cheng, J., Liu, Y., Shi, Q., ... Hao, J. (2019). Drivers of improved PM2.5 air quality in China from 2013 to 2017. \textit{Proceedings of the National Academy of Sciences}, 116(49), 24463--24469.

% ---- References added for the npjUS reframe (VERIFY final bibliographic details before submission) ----
\bibitem[World Bank(2024c)]{worldbank2024tashkent}
World Bank. (2024c). \textit{Air Quality Assessment for Tashkent and the Roadmap for Air Quality Management Improvement in Uzbekistan}. World Bank, Washington, DC.

\bibitem[Malik et al.(2024)]{malik2024classroom}
Malik, S., et al. (2024). Clean air in the classroom: environmental inputs and human capital formation. Working paper (AEA RCT Registry AEARCTR-0012772).

\bibitem[Marshall et al.(2024)]{marshall2024testscores}
Marshall, J. D., et al. (2024). Reducing indoor particulate air pollution improves student test scores: a randomized double-blind crossover study. \textit{Environmental Science \& Technology}.

\bibitem[Banholzer et al.(2024)]{banholzer2024aircleaners}
Banholzer, N., Jent, P., Bittel, P., et al. (2024). Air cleaners and respiratory infections in schools: a modeling study using epidemiological, environmental, and molecular data. \textit{The Lancet Regional Health -- Europe}.

\bibitem[Vieira et al.(2025)]{vieira2025hifive}
Vieira, V., et al. (2025). Classroom air quality in a randomized crossover trial with portable HEPA air cleaners (HIFIVE). California Air Resources Board.

\end{thebibliography}

\end{document}
